% TOP_PROBLEM: {"id": 3070, "title": "The large sets conjecture", "category_id": 2, "category": {"id": 2, "name": "combinatorics", "display_name": "Combinatorics", "description": "Counting problems, graph theory, discrete structures.", "slug": "combinatorics", "order_index": 2, "created_at": "2026-07-31T15:26:25.671Z"}, "status": "open", "problem_number": "OPG-373", "set_id": 12}
% FIRST_POSTED: 2026-10-01
% ENTRY_KIND: statement_only
% UnsolvedMath ID 3070; source catalogue code OPG-373
% !TeX program = lualatex
% Definition and sources only; this file is not an LLM solution attempt.
\documentclass[11pt,a4paper]{article}
\usepackage[margin=25mm]{geometry}
\usepackage{fontspec}
\setmainfont{lmroman10-regular.otf}[BoldFont=lmroman10-bold.otf,
 ItalicFont=lmroman10-italic.otf,BoldItalicFont=lmroman10-bolditalic.otf]
\usepackage{amsmath,amssymb,mathtools}
\usepackage{unicode-math}
\setmathfont{latinmodern-math.otf}
\AtBeginDocument{\let\setminus\smallsetminus}
\usepackage{xurl,hyperref}
\hypersetup{colorlinks=true,urlcolor=blue}
\setcounter{secnumdepth}{0}
\setlength{\parindent}{0pt}
\setlength{\parskip}{5pt}
\setlength{\emergencystretch}{3em}
\begin{document}
\section*{The large sets conjecture}\label{3070}
{\small UnsolvedMath ID 3070 · OPG-373 · Combinatorics · OpenGarden}
\subsection{Definitions and mathematical statement}
Let $L\subseteq\mathbb N=\{1,2,\ldots\}$. For a positive integer $r$,
an $r$-colouring is any map $\chi:\mathbb N\to\{1,\ldots,r\}$; using
all the colours is not required. A monochromatic arithmetic progression
of length $k\ge2$ with an allowed difference is a sequence
\[
 a,a+d,\ldots,a+(k-1)d,\qquad a\ge1,\quad d\in L,
\]
on which $\chi$ takes a single value. Since $L$ contains positive integers,
these terms are distinct and in increasing order.

The set $L$ is $r$-large if every $r$-colouring has such a progression
for every finite length $k\ge2$. It is large if it is $r$-large for every
positive integer $r$.
The conjecture is
\[
 \forall L\subseteq\mathbb N,\qquad
 L\text{ is }2\text{-large}\ \Longrightarrow\ L\text{ is large}.
\]
Consequently the two-colour hypothesis must force monochromatic
progressions of arbitrary finite length in every finite number of colours,
while retaining the same set $L$ of permitted common differences.

The initial term and difference may depend on both the colouring and the
requested length. Neither a fixed difference that works for every length
nor a single infinite progression is demanded. It is essential that each
individual progression uses one common difference throughout: allowing
successive gaps to vary within $L$ gives a different accessibility problem.
No density, additive closure, decidability, or algebraic structure of $L$
is assumed. A counterexample would be one two-large set failing to be
$r$-large for some finite $r>2$.
\subsection{Short English statement}
If allowed differences force arbitrarily long monochromatic progressions in two colours, must they do so in every finite number of colours?
\subsection{Sources}
\begin{itemize}
\item[S1] ULAM.AI and the UnsolvedMath Contributors, supplied problems.json, record 3070 (OPG-373), OpenGarden collection. Catalogue identity and source transcription; CC BY 4.0.
\url{https://huggingface.co/datasets/ulamai/UnsolvedMath}.
\item[S2] Open Problem Garden, The large sets conjecture. Original problem and discussion; the mathematical formulation below is expanded from the supplied source record.
\url{http://www.openproblemgarden.org/op/the_large_sets_conjecture}.
\end{itemize}
\end{document}
